\section{Introduction}

Life on Earth has evolved within the unyielding context of a constant gravitational field. Gravity acts not merely as a mechanical load but as a ubiquitous morphogenetic cue. A 70 kg adult exerts approximately 686 N of axial force through the spine in upright posture. Maintaining this non-geodesic, load-bearing geometry requires continuous metabolic expenditure—an estimated 100 W just for upright standing. This presents a classical biological question: how does genetic information specify and maintain a three-dimensional geometry that actively resists physical forces?

Current theories explain aspects of this process but face critical limitations. Morphoelastic rod theory models how growth tensors prescribe rest curvature but does not directly link to developmental genetics or metabolic energy balance~\cite{goriely2017mathematics}. Biomechanical scoliosis models, such as the Hueter-Volkmann effect, describe a vicious cycle of asymmetric loading and growth modulation \textit{after} a deformity initiates but cannot explain the etiologic origin~\cite{stokes2006hueter}. Hormonal theories based on estrogen or melatonin offer epidemiological correlations but lack a unifying mechanical mechanism. Proprioceptive control theories, such as those implicating PIEZO2, explain alignment maintenance but do not provide a quantitative framework for adolescent vulnerability~\cite{assaraf2020piezo2}. A unifying quantitative framework connecting genetic information, tissue mechanics, metabolic energy, and clinical deformity does not yet exist.

To resolve this, we present the Information-Elasticity Coupling (IEC) framework. We show that coupling a developmental information field $I(s)$ to Cosserat rod rest curvature and stiffness selects the characteristic spinal S-curve as an energy minimum against gravity. While our framework connects developmental patterning to macroscopic geometry, the mechanistic basis lies in protein-level biophysics. HOX genes regulate the synthesis of mechanosensory proteins (e.g., integrins) and structural proteins (e.g., collagens, proteoglycans), which compete for limited metabolic resources during growth. This competition creates thermodynamic instability windows where the system becomes sensitive to perturbations. Furthermore, information-elasticity coupling is fundamentally tensorial: tissue response depends on the alignment between developmental gradient directions and mechanical stress vectors. These protein-level and geometric details provide the mechanistic foundation for IEC.

This approach reveals that the S-shape is maintained at a thermodynamic cost. AlphaFold structural analysis of 23 proteins reveals a 34\% structural cost asymmetry (Demand-Supply Anisotropy Gap) between error-correction and energy-supply subsystems. This creates a predictable Energy Deficit Window during adolescent growth, characterized by six distinct molecular mechanisms of energy supply failure beyond caloric insufficiency. The framework explains the irreversible progression of adolescent idiopathic scoliosis (AIS) via a VIM collapse cascade and generates 12 quantitative, testable predictions.

AIS affects approximately 3\% of adolescents worldwide. Current treatments rely heavily on bracing and invasive fusion surgery because the fundamental etiology remains unknown. The IEC framework provides the first quantitative etiological theory of AIS, introducing patient-stratification metrics (such as the $\chi_\kappa$ parameter) with the potential to predict scoliotic progression. While this work relies on theoretical and computational modeling rather than direct wet-lab experimentation, the mathematical framework and emergent predictions offer a new paradigm for understanding and clinically managing spinal geometry.
